\noindent \textit{\textcolor{red!60!black}{
That said, I think there is one larger issue outstanding. The authors reference the portfolio returns (stock market responses) to various legal developments in section 3 (page 10). This analysis is insightful but not well-explained. In my opinion, the return analysis deserves a table and more detail on the data and summary statistics, some of which can be provided in the appendix. A figure might help too. Consider time series of (cumulative) returns for a) the market, b) firms with gold exposure (unweighted), firms with gold exposure (Weighted as indicated in footnote 10) and d) firms without gold exposure. Such a figure with ``event lines'' (joint resolution, first lawsuits, potentially lower court rulings, start of Supreme Court arguments (poor government showing), Supreme Court decision) would visualize the argument of the authors. I could also see the return analysis in a regression framework. The point here is that illustrate that market clearly responded to relevant events around the abrogation supporting the key identification argument of the authors.
}}

\bigskip

\noindent Thank you for this suggestion. We have added a table (Table~1) and three event study figures (Figures~3--5) to the paper. Each figure now plots exactly the four series the referee lists: (a) the CRSP value-weighted market, (b) an equal-weighted portfolio of firms with gold-clause exposure, (c) the gold-exposure-weighted portfolio of footnote~10, and (d) an equal-weighted portfolio of firms without gold-clause exposure. All four series are constructed from CRSP daily returns, with firms classified by the paper's exposure measure $d_j$. Table~1 reports the raw returns of all four portfolios over each event window, along with the firm counts; Internet Appendix Table~IA.20 reports CAPM cumulative abnormal returns---with the portfolio betas and estimation methodology detailed in its notes---for the intermediate litigation events, with the Joint Resolution window included for reference.

\medskip
\noindent \textit{Main evidence.} The central result is Event~1: the Joint Resolution (May~26--June~6, 1933). Over this nine-day window the market gained $12.1\%$, while the equal-weighted portfolio of gold-exposed firms gained $30.9\%$ and the exposure-weighted portfolio $31.4\%$. The four-series comparison the referee requested is informative here because it separates two components of these gains. First, non-exposed firms---which are comparable in size to the exposed firms---also beat the market, returning $25.4\%$: the reflation rally was strongest among smaller firms, and both equal-weighted portfolios load on that segment relative to the value-weighted market. Second, and central to our identification argument, \textit{gold-exposed firms outperformed observably similar non-exposed firms}: the exposed-minus-non-exposed differential in CAPM cumulative abnormal returns is $+3.1$ percentage points over nine trading days ($+15.8\%$ vs.\ $+12.7\%$; $t \approx 2.1$ relative to the calm-period variability of the daily differential; Table~IA.20). This within-sample comparison nets out the size component and isolates the gold-specific response that the abrogation's debt relief predicts. Event~3 produces a differential in the right direction but of modest size, while the three-day oral-arguments window is essentially flat (we return to it below), and the main identification evidence rests on Event~1, corroborated by the formal regression analysis in Section~5.

\medskip
\noindent \textit{On the oral arguments event.} The three-day window ($-1.3\%$ market, $-1.4\%$ gold-weighted) understates the full impact visible in Figure~4. Over the six trading days following the arguments (January~11--17), the gold-weighted portfolio declined a further $4.6\%$ against a market decline of $3.0\%$, as news of the government's poor showing filtered through press coverage. By January~13, the \textit{New York Times} reported Liberty Bond prices at their highest level since 1917 \citep{NYT1935}, reflecting investor expectations that the government would lose the case. From January~8 through January~17, the cumulative decline was $6.0\%$ for the gold-weighted portfolio and $4.3\%$ for the market. Section~3 now makes this point directly, so that the reader sees why the three-day window alone understates the market response to the arguments.

\medskip
\noindent \textit{Full-period figure with event lines.} We also constructed the single full-period cumulative-return series with event lines that the referee describes. Over a two-year window dominated by the 1933 reflation rally, however, the cumulative scale compresses the event-window movements that carry the identification, and the figure largely duplicates the information in Figures~3--5. We therefore present the event-line evidence at the event level---Figures~3--5 plot all four series around each key event---and examine the remaining litigation events in Internet Appendix Section~IA.1.

\medskip
\noindent \textit{First lawsuits and lower court rulings.} The referee's parenthetical asks about the first lawsuits and lower court rulings, and we now examine every intermediate legal step of the litigation directly. New Internet Appendix Section~IA.1 (Table~IA.20) applies the same portfolios and CAR methodology to: the first court hearing on the enforceability of gold clauses (May~22--24, 1933); the first federal ruling upholding the Joint Resolution's constitutionality (\textit{In re Missouri Pacific R.~Co.}, E.D.~Mo., June~20, 1934); the New York Court of Appeals' affirmance of \textit{Norman v.\ Baltimore \& Ohio R.~Co.} (July~3, 1934); and the Supreme Court's certiorari grants (October~8 and November~5, 1934). Apart from the November grant---which coincided with the midterm election and is discussed separately below---none of these judicial events produced a significant differential response between gold-exposed and non-exposed firms (differentials of $+0.7$, $-0.9$, $+0.4$, and $+0.6$ percentage points for the first hearing, the two lower-court rulings, and the October grant, respectively, all with $|t| \le 1$). We interpret these null results as evidence that the procedural steps were largely anticipated: every court to consider the question in 1933--34 upheld the Joint Resolution, so these rulings confirmed the status quo rather than conveying news. The differential responses appear precisely at the events that carried genuine constitutional and political news---the Joint Resolution and the Supreme Court endgame---which we view as reinforcing the paper's event selection. A footnote in Section~3 now points readers to this analysis.

\medskip
\noindent \textit{The November 1934 window.} The one intermediate window with a significant differential is November~5--8, 1934 ($+3.1$ percentage points, $t \approx 3.6$), and we discuss it in Section~IA.1 (Figure~IA.2) rather than in the main text for a specific reason: it confounds two events. Certiorari in \textit{United States v.\ Bankers Trust Co.} was granted on Monday, November~5; the midterm election was held on Tuesday, November~6, with the New York Stock Exchange closed; and trading on the election result---a Democratic landslide widely read as a popular ratification of the New Deal, plausibly raising the perceived durability of the abrogation---began on Wednesday, November~7. The differential opens over the window's three trading days (November~5, 7, and 8) and does not subsequently revert. Because the certiorari grant and the election cannot be separated within the window, and because we identified this window in the course of the analysis rather than ex ante, we present it as corroborating evidence---a second, politically driven shock to the abrogation's expected fate moving gold-exposed firms differentially---rather than as a main event.

\medskip
\noindent \textit{On the regression framework.} The regression-based identification evidence is developed formally in Section~5. At the event level, the differential CARs in Table~IA.20---each scaled by its calm-period standard error---provide the regression-style complement the referee suggests: they test, event by event, whether gold-exposed firms responded differently from non-exposed firms.
